\subsection{Coverings}

DEFINITION 2. --- A covering is called a locally trivial fiber space all of whose fibers are discrete.

Instead of saying that a B-space E is a covering, one also says that E is a covering of B.

Let E be a B-space, denote by p its projection; let A be a subset of B. If E is a covering of B, the fibers of p are discrete, hence those of \(p_{\mathrm{A}} \colon \overline{p}^{1}(\mathrm{A}) \to \mathrm{A}\) also, and the locally trivial fibered A-space ( \(\overline{p}^{1}(\mathrm{A}), p_{\mathrm{A}}\) ) is a covering.

PROPOSITION 1. --- Let B and E be topological spaces, and let \(p\colon \mathrm{E}\to \mathrm{B}\) be a map. The following conditions are equivalent:

\begin{enumerate}
\def\labelenumi{(\roman{enumi})}
\item
  The map p is continuous and the B-space (E,p) is a trivializable covering.
\item
  There exists a partition \((\mathrm{V}_{i})_{i\in\mathrm{I}}\) of E consisting of open sets such that, for every \(i\in I\), the map \(p\mid V_{i}\colon V_{i}\to B\) is a homeomorphism.
\end{enumerate}

Suppose that condition (i) is satisfied, and let \(g \colon \mathrm{E} \to \mathrm{B} \times \mathrm{F}\) be a trivialization of the B-space \((\mathrm{E}, p)\), of fiber type F. Consequently, F is a discrete topological space and the sets \(\mathrm{V}_i = g^{-1}(\mathrm{B} \times \{i\})\), for \(i \in F\), form a partition of E consisting of open sets. For every \(i\), the map \(p \mid V_i \colon V_i \to B\) is a homeomorphism, with inverse homeomorphism the map \(x \mapsto g^{-1}(x, i)\), whence condition (ii).

Conversely, suppose that condition (ii) is satisfied. The map p is then continuous. Consider the map from E into \(B \times I\) that sends \(x \in E\) to the pair \((p(x), i)\), where i is the unique element of I such that \(x \in V_i\). If I is equipped with the discrete topology, condition (ii) means that this map is a B-isomorphism and the B-space \((E, p)\) is a trivializable covering.

COROLLARY 1. --- Let B and E be topological spaces and let \(p\colon \mathrm{E}\to \mathrm{B}\) be a map. For E, equipped with the map \(p\), to be a covering of B, it is necessary and sufficient that, for every point \(a\) of B, there exists an open neighborhood U of \(a\) and a partition \((\mathrm{V}_i)_{i\in \mathrm{I}}\) of \(\overline{p}^{1}(\mathrm{U})\) formed by open sets of E such that, for every \(i\in \mathrm{I}\), the map \(p\) induces a homeomorphism of \(\mathrm{V}_i\) onto U.

Remark. --- Let B be a topological space. Let E be a B-space and let p be its projection. If E is a covering of B, it follows from Corollary 1 above that the map p is étale (I, p. 28, def. 2) and separated (I, p. 25, prop. 1, (iii)). These conditions are not, however, sufficient to ensure that E is a covering. Consider, for example, an open set U of B. The canonical injection i: U → B is étale and separated. For the B-space (U, i) to be a covering, it is necessary and sufficient that the open set U be also closed.

COROLLARY 2. --- Let B be a topological space, let E be a B-space and let p be its projection. Suppose that p is étale and separated and that the space B is connected. For E to be a trivializable covering of B, it is necessary and sufficient that, for every point x of E, there exists a continuous section s: B → E of p such that \(s(p(x)) = x\).

This is obviously necessary. Conversely, for any continuous section \(s\) of \(p\), the set \(s(\mathrm{B})\) is open and the map \(p\) induces a homeomorphism of \(s(\mathrm{B})\) onto B (I, p. 30, cor. 3 of prop. 6). Moreover, as \(s\) ranges over the set \(\mathcal{S}(\mathrm{B};\mathrm{E})\) of sections of \(p\), the sets \(s(\mathrm{B})\) are pairwise disjoint, since B is connected (I, p. 34, cor. 1 of prop. 11). It then follows from Proposition 1 that if the sets \(s(\mathrm{B})\) cover E, the space E is a trivializable covering of B.

